% !TeX root = ../main.tex

\chapter{Related Work}\label{chapter:related-work}

\section{Compression and the Placement of Visual Analysis}

Compress-then-analyze and analyze-then-compress allocate different work to a
capturing device and its receiver. Hofer et al.~\cite{hofer2023h264cta}
investigate H.264 transmission for edge-assisted visual SLAM, evaluating
encoding, transmission, decoding, and feature-analysis latency on a Raspberry
Pi and an Intel NUC. Their software OpenH264 experiments motivate transmitting
compressed video when capture-side feature extraction is costly; they evaluate
the coding pipeline independently of a complete SLAM system and do not measure
energy consumption.

Redondi et al.~\cite{redondi2016cta} compare JPEG-based CTA and BRISK-feature ATC
for image recognition in visual sensor networks. Their measurements and network
evaluation relate recognition accuracy, transmission rate, and camera-node
energy. The outcome depends on the hardware, feature workload, and link
constraints. These studies motivate the allocation of visual computation; the
present thesis investigates what the receiver can infer directly from codec
motion without a pixel-based correspondence stage.

\section{Motion-Vector Refinement into Optical Flow}

Young, Girod, and Taubman~\cite{young2020fast} recover a dense optical-flow field from block-based coded motion using confidence-weighted filtering and inpainting. Prediction-unit dimensions inform the confidence model, while decoded image colours guide edge-preserving filtering; optional variational refinement additionally uses image gradients, temporal intensity differences, and warping. This work characterizes how codec motion can initialize a more accurate flow estimate, but its recovered flow is a pixel-guided product rather than a codec-only geometric observation.

\section{Codec-Aware Tracking}

VOCA~\cite{hilscher2026voca} uses H.264 motion vectors as initialization priors
for KLT tracking. The codec prediction moves the initial estimate closer to the
subsequent photometric solution, while KLT refinement and conventional tracking
modes retain access to decoded image intensities. VOCA provides the implementation
context and the reference configurations used in this thesis, but does not
impose the same pixel-access restriction.

\section{Motion Vectors for Camera Motion and Mapping}

Li and Song~\cite{li2014toward} estimate monocular camera motion and planar scene
structure from MPEG motion vectors. Their formulation explicitly assumes a
structured environment dominated by planes, such as building facades and paved
roads. It merges direct and bidirectional-frame observations, rejects
insufficient-parallax motion, groups correspondences with homography RANSAC, and
jointly tracks planes and camera motion in an extended Kalman filter. It thus
establishes an early codec-only navigation result within this planar setting,
rather than a general point-based reconstruction. Its B-frame observations also
depend on subsequent anchor images: they cannot be formed as soon as each
intermediate image is captured, introducing buffering and latency for a live
camera. The paper does not establish end-to-end real-time operation. These
constraints concern both the required scene structure and the availability of
the measurements, independently of the estimator's computational cost.

MoV-SLAM~\cite{turner2023movslam} uses H.264 macroblocks and their motion vectors to build long-range correspondence tracks within a complete monocular and stereo SLAM system. Its EXPRESS detector evaluates block intensities to select distinctive regions, while Lucas-Kanade tracking handles missing motion-vector coverage and stereo association. It consequently demonstrates how codec motion can reduce frontend cost without enforcing the pixel-access boundary studied here.

\section{Motion Vectors for Correspondence and Reconstruction}

Zouein et al.~\cite{zouein2025leveraging} convert AV1 block centres and motion vectors into point correspondences, propagate them into multi-frame tracks, and reject directionally inconsistent tracks using cosine similarity. Their evaluation focuses on correspondence density, pairwise geometry, and a small COLMAP structure-from-motion demonstration rather than a dedicated visual-odometry estimator.

The subsequent Gaussian-splatting work~\cite{zouein2026efficient} reuses this
correspondence pipeline to produce dense COLMAP point clouds for initializing 3D
Gaussian splatting. It evaluates the effect of denser reconstruction on
rendering quality and training time, without introducing calibrated stereo depth
or a motion-vector-specific visual-odometry backend.

\section{Motion-Vector Stereo}

Passthrough+~\cite{chaurasia2020passthrough} repurposes video-encoding hardware
to estimate correspondences between rectified views of a stereo pair. Shifted
copies of the stereo images reduce the displacement presented to the encoder
and offset the codec block grid, after which consistency checks and calibrated
triangulation produce sparse 3D points. This per-frame stereo stage uses codec
motion, but the complete scene proxy is not motion-vector-only: the system
assumes a high-precision headset pose, typically supplied by an independent
SLAM system, and uses pose to propagate the previous depth map into the current
view. That pose-assisted temporal model is unavailable under the information
restriction of this thesis. The present work adapts the sparse
stereo-reconstruction principle; the differences in preprocessing and downstream
use are specified in \autoref{chapter:method}.

\section{Spatial Feature Management}

Feature-based estimators generally regulate both the number and spatial
distribution of their observations. VINS-Mono maintains a target population of
corner tracks and enforces a minimum separation between newly detected features
to obtain a more uniform distribution~\cite{qin2018vins}. Similarly, the
low-latency VIO pipeline of K\"uhne et al. suppresses nearby detector candidates
so that the same local feature is not represented repeatedly~\cite{kuhne2025lowlatency}.
Both systems make these decisions from image-derived features. Their common
design objective motivates the pixel-free analogue developed here: codec
observations that are indistinguishable at the available spatial and depth
resolution should not provide duplicated geometric support.

\section{Positioning of This Work}

The compressed-domain systems above span the distinctions central to this thesis: whether
codec motion is refined using pixels, whether recovered structure is planar or
point-based, whether metric scale is supplied by stereo, and whether
correspondences persist beyond a frame pair. Work that merely accelerates video
decoding or generic optical flow without using codec side information is outside
this comparison.

Prior codec-only navigation establishes that motion vectors can support
camera-motion estimation in structured, plane-dominated
environments~\cite{li2014toward}. Other work uses codec motion to initialize
pixel-domain tracking or provide correspondences for reconstruction. The present
work investigates one way to use these measurements in a calibrated, point-based
stereo visual-odometry system without requiring a planar scene model. The
evaluated design combines metric motion-vector stereo, persistent 3D
association, suppression of geometrically indistinguishable observations, and
displacement-based optimization without pixel refinement. Its reverse-temporal
streams are nevertheless prepared offline, so this study also leaves causal,
end-to-end real-time operation unverified.

\begin{table}[p]
  \centering
  \caption[Comparison of compressed-domain visual methods]{Comparison of directly related compressed-domain systems. ``Pixels''
  refers to image-intensity access in correspondence generation or geometric
  estimation, not to the image input required by an encoder.}
  \label{tab:related-work-comparison}
  \begingroup
  \footnotesize
  \setlength{\tabcolsep}{1mm}
  \begin{tabular}{@{}L{26mm}L{29mm}L{24mm}L{36mm}L{25mm}@{}}
    \toprule
    Work & Input and representation & Scale and pixel access &
    Track persistence and validation & Estimator or output \\
    \midrule
    Young et al.~\cite{young2020fast}
      & Coded MVs\newline Dense 2D flow
      & No scale\newline Pixels used
      & \textit{Persistence:} none\newline \textit{Validation:} pairwise confidence
      & Pixel-guided flow \\
    Li and Song~\cite{li2014toward}
      & MPEG MVs\newline Planes and pose
      & Monocular\newline No pixels
      & \textit{Persistence:} plane tracks\newline \textit{Validation:} homography RANSAC
      & Extended Kalman filter \\
    MoV-SLAM~\cite{turner2023movslam}
      & H.264 MVs and LK\newline Point tracks
      & Mono/stereo\newline Pixels used
      & \textit{Persistence:} point tracks\newline \textit{Validation:} LK fallback
      & SLAM backend \\
    Zouein et al.~\cite{zouein2025leveraging}
      & AV1 MVs\newline 2D point tracks
      & Monocular\newline No pixels
      & \textit{Persistence:} MV tracks\newline \textit{Validation:} directional check
      & Matching and SfM study \\
    Zouein et al.~(2026)~\cite{zouein2026efficient}
      & AV1 MVs\newline Dense 3D points
      & Monocular SfM\newline Pixels downstream
      & \textit{Persistence:} MV tracks\newline \textit{Validation:} geometric checks
      & COLMAP and 3DGS \\
    VOCA~\cite{hilscher2026voca}
      & H.264-seeded KLT\newline Point tracks
      & Calibrated VO\newline Pixels used
      & \textit{Persistence:} KLT tracks\newline \textit{Validation:} fallback modes
      & Sliding window BA \\
    Passthrough+~\cite{chaurasia2020passthrough}
      & Stereo encoder MVs and SLAM pose\newline Sparse depth and dense proxy
      & Metric stereo\newline Pixels used
      & \textit{Persistence:} temporal depth\newline \textit{Validation:} stereo consistency
      & View synthesis \\
    Pixelless VO (this thesis)
      & H.264 temporal and stereo MVs\newline Persistent 3D points
      & Metric stereo\newline No estimator pixels
      & \textit{Persistence:} landmark IDs\newline \textit{Validation:} reciprocal + TEASER++
      & Sliding window BA \\
    \bottomrule
  \end{tabular}
  \endgroup
\end{table}

No row alone defines the contribution. In particular, codec-only navigation
predates this thesis, motion-vector tracks have already been used for
reconstruction, and motion-vector stereo originates in Passthrough+. The present
study tests how these ingredients behave together under a pixel-free estimation
boundary and within a point-based visual-odometry estimator.
